% This file is part of JTS-MM (https://github.com/paulotto/jts-mm).
% SPDX-FileCopyrightText: 2026 Paul-Otto Müller
% SPDX-License-Identifier: CC-BY-4.0

\section{Design rules common to all concepts}\label{sec:principles}
\begin{enumerate}[label=\textbf{R\arabic*}]
\item \textbf{Put the spring on the lingual side of the molars}, pressing below the lingual height
of contour. The buccal jaw is rigid and only reciprocates.
\item \textbf{Keep the reference rigid and the spring separate.} The marker is only as still as
the part it is mounted on. Whatever defines the position (rests, reciprocal jaw, frame) should
not flex; the spring only supplies preload. A preloaded rigid contact does not move until the
disturbing force exceeds the preload.
\item \textbf{Use rests as the vertical datum.} Small pads on the marginal ridges next to the
embrasure between the second premolar and first molar define how far the part sits down, the way
occlusal rests do on a partial denture. They replace the incisal stop of the old design.
\item \textbf{Cross the bite only at embrasures}, and check those points with articulating paper.
Everything else stays off the occlusal surfaces, so chewing kinematics are not changed.
\item \textbf{Block out, then relieve.} On the model, fill undercuts deeper than the target
engagement (\SIrange{0.3}{0.5}{\milli\metre} for polymers, \SI{0.4}{\milli\metre} for the
\SI{0.9}{\milli\metre} arch wire) and add \SIrange{0.1}{0.2}{\milli\metre} relief on every
surface that should not touch.
\item \textbf{Reline with resin, and give it something to hold.} A thin layer of light- or
cold-curing resin captures the real tooth shape, including the interproximal undercuts. Resin does
not bond to PETG, PA12 or TPU: design holes, grooves or undercut slots into the printed part.
\item \textbf{Materials for participants.} FDM PETG, PA12 and TPU are fine for bench prototypes
on a printed model but are not certified for intraoral use. Parts worn by participants should be
printed in a certified splint resin (e.g.\ a Class~IIa occlusal-splint resin; flexible grades
exist as an alternative to TPU) and fitted under dental supervision, as covered by the ethics
approval. Both stainless wire grades in \cref{sec:procurement} contain about \SI{8}{\percent}
nickel: ask about nickel allergy.
\end{enumerate}
